\subsection{无穷小变换}

设 \(r \in \mathbf{N}_{\kappa}\)，且 \(X\) 是 \(\mathbf{C}^{r+1}\) 类流形。

8.4.1. 设 \(\tau\) 是一个 \(C^r\) 类同构向量函子（7.6.6）。若 E 是 Banach 空间，记 \(\mathbf{GL}(\mathbf{E})\) 为由 E 的自同构组成的 \(\mathcal{L}(\mathbf{E};\mathbf{E})\) 的开子集。记 \(\tau_{\mathbf{E}}'\) 为态射 \(\mathrm{Id}_{\mathbf{E}}\) 在单位元 \(\tau\colon\mathbf{GL}(\mathbf{E})\to\mathbf{GL}(\tau(\mathbf{E}))\) 处的线性切映射；它是从 \(\mathcal{L}(\mathbf{E};\mathbf{E})\) 到 \(\mathcal{L}(\tau(\mathbf{E});\tau(\mathbf{E}))\) 的连续线性映射。

8.4.2. 设 E 是 Banach 空间，并令 \(F = \tau(E)\)。设 U 是 E 的开子集，\(\xi\) 是 U 上的 \(C^r\) 类向量场，\(\tilde{\xi} \colon U \to E\) 是相应的映射（通过将 T(U) 与 U \times E 等同而得到（8.1.1）），且 \(f \colon U \to F\) 是一个 \(C^r\) 类映射。对于 \(x \in U\)，按下式定义 F 中的元素 \((\theta_\xi.f)(x)\)：

\[
(\theta_ {\xi}. f) (x) = d _ {x} f (\tilde {\xi} (x)) - \tau_ {\mathrm{E}} ^ {\prime} (d _ {x} \tilde {\xi}) (f (x)).\tag{1}
\]

（注意，一方面，\(\tilde{\xi}(x) \in \mathbf{E}\) 且 \(d_x f \in \mathcal{L}(\mathbf{E}; \mathbf{F})\)，因而 \(d_x f(\xi(x)) \in \mathbf{F}\)；另一方面，\(d_x \tilde{\xi} \in \mathcal{L}(\mathbf{E}; \mathbf{E})\)，\(\tau_{\mathbf{E}}'(d_x \tilde{\xi}) \in \mathcal{L}(\mathbf{F}; \mathbf{F})\)，且 \(f(x) \in \mathbf{F}\)，因而 \(\tau_{\mathbf{E}}'(d_x \tilde{\xi})(f(x)) \in \mathbf{F}\)。函数 \(\theta_{\xi}.f: x \mapsto (\theta_{\xi}.f)(x)\) 在 U 中属于 \(\mathbf{C}^{r-1}\) 类。）

8.4.3. 令 F 表示向量丛 \(\tau(\mathrm{T}(\mathrm{X}))\)（7.6.2 和 7.6.6）。设 \(\xi\) 是 X 上的向量场，\(f\) 是 F 的截面，二者都属于 \(\mathbf{C}^r\) 类。设 \(x\) 是 X 的一点，并设坐标图 \(c = (\mathrm{U}, \varphi, \mathrm{E})\)，它是 X 在 \(x\) 处的坐标图。令 \(c' = (\mathrm{U}, \zeta_c, \mathrm{E})\) 为 T(X) 相应的向量坐标图（8.1.1），令 \(\tau(c') = (\mathrm{U}, \psi_c, \tau(\mathrm{E}))\) 为 F 相应的向量坐标图（7.6.2）。置 \(x_c = \varphi(x)\)；令 \(\xi_c\) 为从 \(\varphi(\mathrm{U})\) 到 E 的映射，使得 \(\zeta_c(\xi(u)) = (u, \xi_c(\varphi(u)))\)，并令 \(f_c\) 为从 \(\varphi(\mathrm{U})\) 到 \(\tau(\mathrm{E})\) 的映射，使得 \(\psi_c(f(u)) = (u, f_c(\varphi(u)))\) 对所有 \(u \in \mathrm{U}\) 成立。将 8.4.2 中定义的 \(a_c\)（即 \((\theta_{\xi_c}.f_c)(x_c)\)）看作 \(\tau(\mathrm{E})\) 中的元素。存在且仅存在一个元素 \(a\)，属于 \(\mathrm{F}_x\)，使得 \(\psi_c(a) = (x, a_c)\) 对 X 的每个坐标图 \(c\)（在 \(x\) 处）都成立。记此元素为 \((\theta_\xi.f)(x)\)。它只取决于 \(\xi\) 和 \(f\) 在 \(x\) 处的芽（甚至只取决于它们的一阶喷射（12.1.2））。

映射 \(\theta_{\xi}.f: x \mapsto (\theta_{\xi}.f)(x)\) 是 F 的 \(\mathbf{C}^{r-1}\) 类截面；它关于对 \((\xi, f)\) 是 K-双线性的。

8.4.4. 保持前述假设和记号。存在（7.6.4）一个向量丛态射 \(\tau'\)；从 \(\mathcal{L}(\mathrm{T}(\mathrm{X});\mathrm{T}(\mathrm{X}))\) 到 \(\mathcal{L}(\mathrm{F};\mathrm{F})\) 的这种态射仅有一个，使得对所有 \(x \in X\)，\(\tau'\) 在纤维 \(\mathcal{L}(T(X); T(X))_x = \mathcal{L}(T_x(X); T_x(X))\) 上的限制等于 \(\tau_{T(X)}'\)（8.4.1）。若 \(g\) 是 \(C^r\) 类函数，定义在 \(X\) 上且取值于 \(K\)，则有

\[
\theta_ {g \xi}. f = g \theta_ {\xi}. f - \tau^ {\prime} (d g \otimes \xi) (f)
\]

（在此公式中，将 \(dg \otimes \xi\) 视为 \(T'(X) \otimes T(X)\) 的截面，并将其视为 \(\mathcal{L}(T(X); T(X))\) 的截面；其像 \(\tau'(dg \otimes \xi)\) 是 \(\mathcal{L}(F; F)\) 的截面，并将 \(f\) 映为 \(F\) 的截面）。

8.4.5. 设 I 是包含 0 的 K 的开子集，U 是 X 的开子集，且 \(\varphi\) 是从 \(C^r\) 到 X 的 \(\mathbf{I} \times \mathbf{U}\) 类映射。对于 \((t, x) \in \mathbf{I} \times \mathbf{U}\)，记向量 \(\varepsilon_{t,x}\) 为 \((1, 0) \in \mathbf{K} \times \mathrm{T}_x(\mathbf{X}) = \mathrm{T}_{(t,x)}(\mathbf{I} \times \mathbf{U})\)。对于 \(t \in \mathbf{I}\)，记从 U 到 X 的映射 \(\varphi_t\)，其定义为 \(\varphi_t(x) = \varphi(t, x) (x \in \mathbf{U})\)。假定满足下列三个条件：

\begin{enumerate}
\def\labelenumi{\alph{enumi})}
\item
  \(\varphi_{0}\) 是从 U 到 X 的典范嵌入；
\item
  对所有 \(t \in \mathbf{I}\)，映射 \(\varphi_t\) 是从 \(\mathbf{C}^r\) 到 \(\mathbf{U}\) 的某个开子集的 \(\mathbf{X}\) 类流形同构；
\item
  丛 \((t, x) \mapsto T_{(t, x)}(\varphi)(\varepsilon_{t,x})\) 的截面 \(\varphi^*T(X)\) 属于 \(C^r\) 类。此外，置：
\end{enumerate}

\[
d) \xi (x) = \mathrm{T} _ {(0, x)} (\varphi) (\varepsilon_ {0, x}) \quad \text {for all} \quad x \in \mathrm{U}.
\]

映射 \(\xi\colon x\mapsto\xi(x)\) 是 U 上的 \(\mathbf{C}^{r}\) 类向量场。称其为 \(\varphi\) 的初始（或起始）场。

设 \(\tau\) 是一个 \(C^r\) 类同构向量函子；置 \(F = \tau(T(X))\)，并设 \(f\) 是 \(C^r\) 类的、取值于 \(F\) 且定义在 \(X\) 上的截面。对于 \(x \in U\) 和 \(t \in I\)，置 \(\varphi_t^* f(x) = \tau(T_x(\varphi_t)^{-1})(f(\varphi_t(x)))\)；它是 \(F_x\) 中的元素。映射 \(t \mapsto \varphi_t^* f(x)\) 是从 \(I\) 到 \(F_x\) 的 \(C^{r-1}\) 类映射，对所有 \(x \in U\) 都成立，且有

\[
\frac {d}{d t} \left(\varphi_ {t} ^ {*} f (x)\right) _ {t = 0} = \left(\theta_ {\xi}. f\right) (x) \quad \text {for all} \quad x \in \mathrm{U}.\tag{2}
\]

8.4.6. 给定 X 上的 \(\xi\) 类向量场 \(C^r\) 及 X 的一点 \(x_0\)，存在包含 0 的 K 的开区间 I、包含 \(x_0\) 的 X 的开集 U，以及一个 \(\varphi: I \times U \to X\) 类映射 \(C^r\)，满足 8.4.5 的条件 \(a\)、\(b)\) 和 \(c)\)，且 \(\varphi\) 的初始场是 \(\xi|U\) 在 U 上的限制 \(\xi\)。

8.4.7. 例。--- 设 F 是 Banach 空间，p 是一个 \(\geqslant0\) 的整数。若 V 是 Banach 空间，置 \(\alpha_{p}(\mathbf{V}) = \operatorname{Alt}^{p}(\mathbf{V}; \mathbf{F})\)；若 u 是 V 到 Banach 空间 \(V'\) 的同构，记由 u 通过结构传递得到的从 \(\alpha_{p}(u)\) 到 \(\alpha_{p}(\mathbf{V})\) 的同构为 \(\alpha_{p}(\mathbf{V}')\)。由此得到一个同构向量函子，记为 \(\alpha_{p}\)（7.8.1 及第 7.8 号勘误）。可将前述结果应用于它；有

\[
\alpha_ {p} (\mathrm{T} (\mathrm{X})) = \operatorname{Alt}^{p}(\mathrm{T} (\mathrm{X}); \mathrm{F}).
\]

\(\alpha_{p}(\mathrm{T}(\mathrm{X}))\) 的截面是 \(p\) 次、定义在 \(\mathbf{X}\) 上且取值于 \(\mathbf{F}\) 的微分形式；若 \(\omega\) 是这样的形式，且 \(\xi\) 是 \(\mathbf{X}\) 上的向量场，二者都属于 \(\mathbf{C}^{r}\) 类，则 \(\theta_{\xi}.\omega\) 是 \(p\) 次、定义在 \(\mathbf{X}\) 上、取值于 \(\mathbf{F}\) 且属于 \(\mathbf{C}^{r-1}\) 类的微分形式。有

\[
\theta_ {\xi}. \omega = d (i (\xi) \omega) + i (\xi) d \omega\tag{3}
\]

并且当 \(r \geqslant 2\) 时，

\[
\theta_ {\xi}. d \omega = d (\theta_ {\xi}. \omega).\tag{4}
\]

当 \(p=0\) 时，函子 \(\alpha_{p}\) 是由 F 定义的常值向量函子；丛 \(\alpha_{0}(\mathrm{T}(\mathrm{X}))\) 与平凡丛 \(\mathrm{F}_{\mathrm{x}}\) 等同，而该丛的截面与 X 上取值于 F 的函数等同。若 \(\xi\) 是 X 上的 \(\mathbf{C}^{r}\) 类向量场，且 \(f\in\mathcal{C}^{r}(\mathrm{X};\mathrm{F})\)，则有

\begin{enumerate}
\def\labelenumi{(\arabic{enumi})}
\setcounter{enumi}{4}
\tightlist
\item
  \[
  \theta_ {\xi}. f = \langle \xi , d f \rangle = D _ {\xi} (f) \quad \text{(8.2.3)};
  \]
\end{enumerate}

它是 \(\mathcal{C}^{r-1}(X;F)\) 中的元素。

8.4.8. 设 \(\tau\) 和 \(\tau_{1}\) 是两个同构向量函子。向量函子的一个态射 \(h\colon \tau_{1}\to \tau\) 定义了从 \(h\) 到 \(\tau_{1}(\mathrm{T}(\mathbf{X}))\) 的向量丛态射（仍记为 \(\tau (\mathrm{T}(\mathbf{X}))\)）（7.6.4），并将 \(f\) 的一个截面 \(\tau_{1}(\mathrm{T}(\mathbf{X}))\) 对应为 \(h(f)\) 的一个截面 \(\tau (\mathrm{T}(\mathbf{X}))\)。若 \(\xi\) 是 \(\mathbf{X}\) 上的向量场，且 \(f\) 和 \(\xi\) 属于 \(\mathbf{C}^r\) 类，则有 \(\theta_{\xi}.h(f) = h(\theta_{\xi}.f)\)。

更一般地，设 \(\tau\)、\(\tau_{1},\ldots,\tau_{d}\) 是同构向量函子，且 \(h\colon(\tau_{1},\ldots,\tau_{d})\to\tau\) 是一个 \(d\)-线性态射（7.6.4）。设 \(f_{i}\in\mathcal{S}_{\tau_{i}(\mathrm{T}(\mathrm{X}))}^{r}(\mathrm{X})\)（对 \(i=1,\ldots,d\)），且 \(\xi\) 是 X 上的 \(\mathbf{C}^{r}\) 类向量场。有

\[
\theta_ {\xi}. h (f _ {1}, \dots , f _ {d}) = \sum_ {1 \leqslant i \leqslant d} h (f _ {1}, \dots , f _ {i - 1}, \theta_ {\xi}. f _ {i}, f _ {i + 1}, \dots , f _ {d}).\tag{6}
\]

例如，若 \(f\) 是 \(\tau(\mathrm{T}(\mathrm{X}))\) 的截面，\(g\) 是 \(\mathrm{X}\) 上取值于 \(\mathrm{K}\) 的函数，且二者都属于 \(\mathbf{C}^r\) 类，则有

\[
\theta_ {\xi}. (g f) = (\theta_ {\xi}. g) f + g (\theta_ {\xi}. f).\tag{7}
\]

若 \(\omega_{1}\) 和 \(\omega_{2}\) 是两个分别取值于 Banach 空间 \(C^{r}\) 和 \(F_{1}\) 的 \(F_{2}\) 类微分形式，则有

\[
\theta_ {\xi}. (\omega_ {1} \wedge \omega_ {2}) = (\theta_ {\xi}. \omega_ {1}) \wedge \omega_ {2} + \omega_ {1} \wedge (\theta_ {\xi}. \omega_ {2}),\tag{8}
\]

这里的外积是相对于从 \(F_{1} \times F_{2}\) 到 Banach 空间 F 的连续双线性映射而言的。

8.4.9. 设 \(\tau\) 是反变的（8.2.7）、\(\mathbf{C}^r\) 类同构向量函子，\(\varphi: X \to Y\) 是 \(\mathbf{C}^{r+1}\) 类流形的态射，\(\xi\)（分别为 \(\eta\)）是定义在 \(\mathbf{C}^r\)（分别为 \(X\)）上的 \(Y\) 类向量场，且 \(\xi\) 通过 \(\varphi\) 与 \(\eta\) 相关（8.2.6）。对于每个 \(f\)（它是 \(\tau(T(Y))\) 的 \(\mathbf{C}^r\) 类截面），有

\[
\varphi^ {*} (\theta_ {\eta}. f) = \theta_ {\xi}. \varphi^ {*} f.
\]

